\section{Largest Palindrome Product --- Problem 4}
\label{prob:4}
\noindent\textbf{Source:} \url{https://github.com/KO0I/projectEuler/blob/reorganize/solutions-by-language/solutions/c/largest_palindrome_product_4.c}

\noindent\textbf{Assessment:} Dependent on generated input.

\subsection*{Mathematical goal}
Find the maximum \(ab\) for three-digit integers \(a,b\) such that the decimal digits of \(ab\) form a palindrome.

\subsection*{How the C code works}
The program opens \verb|palindrome|, reads each candidate number into \verb|temp|, and tests \verb|a| from 100 to 999. If \verb|temp%a==0|, it tests \verb|b| in the same range and checks \verb|a*b==temp|. A successful product greater than the running \verb|max| replaces it. The CPU time is measured using \verb|clock|.

\subsection*{Correctness and limitations}
The source assumes the input file contains a complete, correctly separated list of palindromes. The historical palindrome generator does not reliably separate output records, so the data path must be verified. The calculated quotient \verb|q=temp/a| is unused; checking whether \verb|q| is a three-digit integer would remove the inner loop. The expected maximum is \texttt{...6609}, but it is not verified by the existing file workflow.


\subsection*{Correction to historical C --- Problem 4}
{\color{red}\textbf{Correction:} The new implementation enumerates products of three-digit integers and tests each product as a palindrome. It no longer relies on the historical palindrome input file or an unused second factor search.}
